\begin{figure}[H]
  \centering
  \begin{tikzpicture}[
    >=Latex,
    main node/.style={circle,draw,minimum size=6.5mm,inner sep=0pt},
    edge label/.style={draw=none,fill=white,inner sep=1pt,font=\footnotesize},
    every edge/.style={draw,->,line width=0.6pt}
  ]
    % Original Graph G
    \begin{scope}[xshift=0cm, yshift=0cm]
      \node[main node] (U)  at (0, 0) {u};
      \node[main node] (X)  at (0, -2.0) {x};
      \node[main node] (V)  at (2.5, -1.0) {v};
      \node[main node] (W1) at (5.0, 0.5) {w$_1$};
      \node[main node] (W2) at (5.0, -2.5) {w$_2$};
      
      \path (U) edge node[edge label, above=2pt, sloped] {$e_1$} (V)
            (X) edge node[edge label, below=2pt, sloped] {$e_4$} (V)
            (V) edge node[edge label, above=2pt, sloped] {$e_2$} (W1)
            (V) edge node[edge label, below=2pt, sloped] {$e_3$} (W2);
            
      \node[draw=none] at (2.5, -3.5) {\textbf{原图:} $G$};
    \end{scope}
    
    % Arrow indicating transformation aligned vertically with the center (-1.0)
    \node[draw=none] (arrow) at (6.6, -1.0) {\huge $\Rightarrow$};
    
    % Line Graph L(G)
    \begin{scope}[xshift=8.5cm, yshift=0cm]
      \node[main node] (E1) at (0, 0) {$e_1$};
      \node[main node] (E4) at (0, -2.0) {$e_4$};
      \node[main node] (E2) at (3.0, 0.5) {$e_2$};
      \node[main node] (E3) at (3.0, -2.5) {$e_3$};
      
      % Shift positions of diagonal labels so they don't visually intersect
      \path (E1) edge node[edge label, above=1pt, sloped] {$(e_1, e_2)$} (E2)
            (E1) edge node[edge label, pos=0.74, above=1pt, sloped] {$(e_1, e_3)$} (E3)
            (E4) edge node[edge label, pos=0.26, above=1pt, sloped] {$(e_4, e_2)$} (E2)
            (E4) edge node[edge label, below=1pt, sloped] {$(e_4, e_3)$} (E3);
      
      \node[draw=none] at (1.5, -3.5) {\textbf{线图:} $L(G)$};
    \end{scope}

  \end{tikzpicture}
  \caption{有向图网络拓扑映射至线图的转换展开结构}
  \label{fig:line_graph_example}
\end{figure}
